\chapter{Catálogo cerrado y protocolo experimental}
\label{anexo:catalogo}

\section{Por qué este anexo es probatorio}

La afirmación central del diseño experimental --- que ningún parámetro se optimizó fuera de un
catálogo fijado de antemano --- sólo puede comprobarse si el catálogo está publicado. Este anexo lo
reproduce íntegro. Cualquier lector puede verificar dos cosas: que el valor ganador de cada variable
pertenece a la rejilla declarada para ella, y que la configuración final no contiene ninguna
variable ajena al catálogo.

\section{Regla de selección}

Un candidato desplaza al incumbente sólo si supera, en este orden, tres filtros.

\begin{enumerate}[leftmargin=*,itemsep=0.35em]
  \item \textbf{Elegibilidad.} Debe tener observaciones suficientes y no caer por debajo de $-0{,}02$
  de Rank-IC en ninguna era disponible. Una configuración que fracase por completo en un régimen
  queda excluida aunque su promedio sea favorable.
  \item \textbf{Puerta pareada.} Debe cumplir una de dos condiciones: \emph{dominar} al incumbente
  --- ventaja media positiva y mejor en más de la mitad de las cohortes --- o ser \emph{no inferior},
  con límite inferior del \textit{bootstrap} pareado al 90\,\% por encima de $-0{,}01$. Las
  diferencias se emparejan por periodo mensual, de modo que dos configuraciones con retardos de
  ejecución distintos puedan compararse sin confundir el desplazamiento de sus fechas con ausencia
  de evidencia.
  \item \textbf{Desempate.} Entre candidatos elegibles se ordena por ventaja pareada. Una diferencia
  inferior a 0,002 se declara empate técnico y se resuelve por simplicidad, no por magnitud. Si
  persiste el empate se recurre, en orden, al Rank-IC medio, a la fracción de cohortes positivas, a
  la variabilidad, a la complejidad y, finalmente, a conservar el incumbente.
\end{enumerate}

La última cláusula importa: ante evidencia indistinguible, el protocolo \emph{no cambia}. Es lo que
explica que dieciséis de las dieciocho decisiones de la pasada de referencia terminen confirmando
el valor de partida, según la traza del Capítulo~\ref{chap:diseno-experimental}. La relación
completa de candidatos, puertas de elegibilidad, ventajas pareadas e intervalos está en
\artefacto{decisions.json}: quince decisiones se resolvieron por \texttt{robust\_rank\_ic} y dos por
\texttt{tie\_simplicity}.

Diecisiete decisiones no cubren las veintiuna variables predictivas del catálogo, y la diferencia
conviene declararla. Cuatro --- \artefacto{model\_family}, \artefacto{objective},
\artefacto{meta\_history\_quarters} y \artefacto{meta\_recency\_weighting} --- se ejecutaron en modo
\texttt{fixed} en la pasada de referencia, es decir, con un único valor alcanzable de su rejilla, de
modo que no había comparación que hacer. Sus valores finales son los que la tabla recoge, pero
provienen de una decisión declarada de antemano y no de una medición pareada, exactamente igual que
los dos empates técnicos.

Las variables de cartera no participan en ninguno de estos filtros. Se ejecutan después de congelar
el ganador y se reportan como diagnóstico, según se detalla en la
Sección~\ref{sec:cartera-rejilla}.

\section{El catálogo completo}

La Tabla~\ref{tab:catalogo} reproduce las treinta y tres variables con su etapa, su rejilla de
valores admisibles y el valor que adoptó el ganador. Va apaisada porque la columna de valores es la
que hace auditable el catálogo: sin ella no podría comprobarse que una alternativa estaba disponible
y no se eligió.

\begin{landscape}
\begin{longtable}{lllp{0.34\linewidth}l}
\caption{Catálogo cerrado completo: las 33 variables, su rejilla y el valor del ganador.}\label{tab:catalogo}\\
\toprule
Variable & Etapa & Predictiva & Valores del catálogo & Ganador \\
\midrule
\endfirsthead
\toprule
Variable & Etapa & Predictiva & Valores del catálogo & Ganador \\
\midrule
\endhead
\midrule
\multicolumn{5}{r}{\footnotesize Continúa en la página siguiente}\\
\endfoot
\bottomrule
\endlastfoot
feature\_preset & Representación & Sí & core, all & all \\
meta\_method & Meta-agente & Sí & equal, stacked rolling free, stacked rolling bounded & stacked rolling bounded \\
model\_family & Modelo & Sí & lightgbm, elastic net & lightgbm \\
snapshot\_step\_months & Temporal & Sí & 1, 3, 6, 12 & 1 \\
target\_size & Cartera informativa & No & 8, 12, 16, 25, 50 & 12 \\
exit\_expected\_alpha\_bps & Cartera informativa & No & 0.0, 100.0, 250.0 & 100.0 \\
fundamental\_momentum & Representación & Sí & False, True & True \\
lgbm\_max\_depth & Modelo & Sí & 3, 4, 6 & 3 \\
meta\_history\_quarters & Meta-agente & Sí & 8, 16 & 16 \\
target\_horizon\_months & Temporal & Sí & 3, 6, 12 & 12 \\
lgbm\_n\_estimators & Modelo & Sí & 100, 200, 400 & 100 \\
market\_regime\_feature & Representación & Sí & False, True & False \\
meta\_recency\_weighting & Meta-agente & Sí & off, linear, exponential & off \\
rotation\_edge\_bps & Cartera informativa & No & 25.0, 50.0, 100.0 & 50.0 \\
train\_lookback\_years & Temporal & Sí & 4, 8, 12 & 8 \\
cash\_policy & Cartera informativa & No & fully invested, opportunity cash & opportunity cash \\
max\_cash\_weight & Cartera informativa & No & 0.0, 0.1, 0.25 & 0.25 \\
minimum\_holding\_period & Cartera informativa & No & none, quarter horizon, half horizon, full horizon & none \\
execution\_lag\_days & Temporal & Sí & 30, 45, 60 & 60 \\
lgbm\_learning\_rate & Modelo & Sí & 0.03, 0.05, 0.1 & 0.03 \\
neutralize\_by\_sector & Representación & Sí & False, True & False \\
rebalance\_drift\_tolerance & Cartera informativa & No & 0.0, 0.1, 0.25 & 0.25 \\
lgbm\_min\_child\_samples & Modelo & Sí & 20, 50, 100 & 50 \\
price\_only\_strictness\_multiplier & Cartera informativa & No & 1.0, 1.5, 2.0 & 1.5 \\
recency\_weighting & Temporal & Sí & off, linear, exponential & off \\
winsorization & Representación & Sí & 0.0, 0.01, 0.025 & 0.0 \\
price\_only\_sell\_only & Cartera informativa & No & False, True & False \\
max\_features\_per\_agent & Representación & Sí & 8, 12, 20 & 12 \\
objective & Temporal & Sí & rank regression, ranking & rank regression \\
sizing\_mode & Cartera informativa & No & equal, alpha proportional & alpha proportional \\
commission\_bps & Cartera informativa & No & 0.0, 5.0, 10.0 & 5.0 \\
feature\_weighting\_mode & Representación & Sí & model native, oos stability prune & oos stability prune \\
slippage\_bps & Cartera informativa & No & 5.0, 10.0, 20.0 & 10.0 \\
\end{longtable}

\end{landscape}

Las cuatro primeras etapas --- temporal, representación, modelo y meta --- son predictivas y pueden
modificar al ganador. La quinta, cartera, no. El coste declarado de cada variable indica qué debe
recalcularse al cambiarla: un ajuste completo, únicamente la recombinación del meta-agente o sólo el
\textit{backtest}. Ese coste es lo que hace viable recorrer el catálogo de forma secuencial en lugar
de cartesiana.

\paragraph{Qué ganador recoge la última columna.} La columna «Ganador» es la del \textbf{Model
Study 3}, incluidas sus variables de cartera, que en esa fase conservaban el valor por defecto del
catálogo. No es, por tanto, la cartera que el trabajo adopta: el Portfolio Study reoptimizó después
seis de las doce variables de cartera y cambió tres de ellas ---
\artefacto{max\_cash\_weight} de 0,25 a 0,0, \artefacto{minimum\_holding\_period} de
\texttt{none} a \texttt{half\_horizon} y \artefacto{rebalance\_drift\_tolerance} de 0,25 a 0,1 ---,
mientras que \artefacto{target\_size}, \artefacto{sizing\_mode} y
\artefacto{coverage\_percentile\_floor} coincidieron con los de partida. La configuración adoptada
se presenta en la Sección~\ref{sec:cartera-ganadora}. Se tabula aquí el ganador del Model Study
porque este anexo documenta el catálogo del protocolo predictivo, no la gestión de cartera.

\paragraph{Por qué seis de las doce.} Las otras seis variables de cartera ---
\artefacto{commission\_bps}, \artefacto{slippage\_bps}, \artefacto{exit\_expected\_alpha\_bps},
\artefacto{rotation\_edge\_bps}, \artefacto{price\_only\_strictness\_multiplier} y
\artefacto{price\_only\_sell\_only} --- declaran rejilla pero no entran en el cartesiano, y la razón
es que no describen una decisión del gestor sino el entorno en el que opera. Las dos primeras son
condiciones de mercado: optimizarlas equivaldría a elegir cuánto se quiere pagar por operar, que no
es una elección disponible. Las cuatro restantes gobiernan umbrales de la mecánica de rotación cuyo
efecto ya está mediado por las seis que sí se recorren. Fijarlas es una decisión declarada de
antemano, no una medición, y como tal acota lo que puede afirmarse: el barrido responde a qué
cartera es mejor \emph{dado} ese entorno, no a cuál sería la mejor sin restricciones.

La diferencia de magnitud entre ambos recorridos justifica por sí sola la elección. El producto
cartesiano de las treinta y tres rejillas asciende a $1{,}9 \times 10^{14}$ combinaciones; incluso
restringido a las veintiuna variables predictivas son $5{,}4 \times 10^{8}$. El recorrido secuencial
evaluó 46 configuraciones en la pasada de referencia. Esa reducción no es sólo computacional: cada
configuración adicional ensayada incrementa la probabilidad de encontrar una ganadora por azar, y es
exactamente lo que el Deflated Sharpe del Capítulo~\ref{chap:resultados} penaliza. Un barrido
cartesiano sobre las variables predictivas habría hecho imposible cualquier afirmación estadística
sobre el resultado.

Ese argumento vale para el plano predictivo y no se extiende al de cartera, donde las seis variables
generan 1.440 combinaciones: cinco órdenes de magnitud menos, y sobre una señal que ya no puede
cambiar. El razonamiento completo está en la Sección~\ref{sec:seleccion-secuencial}.
